\chapter{ХЭРЭГЖҮҮЛЭЛТИЙН ЯВЦ БА ЦААШДЫН ТӨЛӨВЛӨГӨӨ}

\section{Одоогийн хэрэгжилтийн төлөв}

Энэхүү төслийг завсрын шатанд үнэлэхэд бүрэн бүтээгдэхүүн биш, харин ажиллаж эхэлсэн прототип гэж дүгнэх нь зөв. Хүснэгт~\ref{tab:progress}--д одоогийн төлөвийг нэгтгэн үзүүлэв.

\begin{table}[H]
  \centering
  \caption{Одоогийн хэрэгжилтийн төлөв}
  \label{tab:progress}
  \begin{tabularx}{\textwidth}{>{\raggedright\arraybackslash}p{4.6cm} >{\raggedright\arraybackslash}p{2.8cm} X}
    \toprule
    Ажлын хэсэг & Төлөв & Тайлбар \\
    \midrule
    Туршилтын вэб интерфейс & Хийсэн & Прототип турших зориулалттай суурь интерфейс бэлэн боловч эцсийн хэрэглээний үндсэн клиентээр төлөвлөөгүй \\
    AI серверийн суурь endpoint & Хийсэн & Elysia дээр `/ai` endpoint ажиллаж, загвараас стрийм хариу авч байна \\
    SISI MCP сервер & Хийсэн & Локал MCP сервер, health болон session bootstrap endpoint-уудтай \\
    Session баталгаажуулалт & Хийсэн & Cookie/header дээр тулгуурласан session verify, refresh, clear логик нэмэгдсэн \\
    Read-only MCP домэйнүүд & Хийсэн & Overview, schedule, academic, financial, notifications, resources, top menus, e-journal домэйнүүд тодорхойлогдсон \\
    AI ба MCP tool orchestration & Хийгдэж байгаа & AI серверийн message урсгал болон MCP хэрэгслийн бодит дуудагдалтыг нэгтгэх ажил үлдсэн \\
    Chrome extension-ийн интеграц & Хийгдэж байгаа & Extension-оор access token авч, дахин нэвтрэлтгүй session bootstrap хийх үндсэн урсгал боловсруулж байна \\
    Ерөнхий MCP өргөтгөлүүд & Төлөвлөсөн & Календарь, баримт бичиг, экспорт зэрэг төрөл бүрийн нэмэлт интеграц холбох боломжтой болгоно \\
    Багшийн порталын API интеграц & Төлөвлөсөн & Багшийн талын өгөгдөл болон mutation workflow-г MCP хэлбэрт оруулах \\
    Excel-ээс дүн оруулах урсгал & Төлөвлөсөн & Олон оюутны дүнг Excel файлаас системд оруулах ажиллагааг хөнгөвчлөх \\
    CSV export, mutation баталгаажуулалт & Төлөвлөсөн & Аюулгүй байдлын нэмэлт давхарга шаардлагатай \\
    Туршилт ба хэрэглэгчийн үнэлгээ & Төлөвлөсөн & Семестрийн сүүлчийн шатанд хийх \\
    \bottomrule
  \end{tabularx}
\end{table}

\section{Завсрын шатны гол үр дүн}

Одоогийн байдлаар дараах бодит ахиц гарсан. SISI API-д холбогдох MCP архитектурын эхний хувилбар ажиллаж эхэлсэн, оюутны өгөгдлийн хэд хэдэн чухал домэйныг read-only хэрэгсэл хэлбэрт оруулсан, туршилтын интерфейс, AI endpoint, Chrome extension-ийн суурь модуль бүрдсэн, цаашдын интеграцад ашиглах боломжтой модульчлагдсан monorepo бүтэц тогтсон, мөн хэрэгжсэн зүйл болон төлөвлөсөн зүйлсийг ялгаж баримтжуулах суурь бүрдсэн.

\section{Үлдсэн ажлууд}

Дараагийн шатанд AI серверээс MCP хэрэгслүүдийг бодитой дуудах orchestration урсгалыг нэгтгэх, хэрэглэгчийн зорилгоос хамаарсан tool routing болон confirmation flow загварчлах, Chrome extension-оос access token авч session reuse хийх урсгалыг тогтворжуулах, ерөнхий MCP өргөтгөлүүдийг турших, жишээ интеграцын оронд өргөтгөх архитектур бүрдүүлэх, багшийн порталын API-д холбогдох MCP давхаргыг төлөвлөх, Excel файлаас дүн оруулах workflow-г аюулгүй хэрэгжүүлэх шийдлийг загварчлах, CSV export, file output, audit trail логикийг тодорхойлох, мөн хэрэглэгчийн туршилт, жишиг сценари, үнэлгээний шалгуур боловсруулах зэрэг ажлуудыг тэргүүн ээлжид хийх шаардлагатай байна.

\section{Ажлын төлөвлөгөө}

\begin{table}[H]
  \centering
  \caption{Дараагийн шатны ажлын төлөвлөгөө}
  \label{tab:roadmap}
  \begin{tabularx}{\textwidth}{>{\raggedright\arraybackslash}p{3cm} >{\raggedright\arraybackslash}X}
    \toprule
    Хугацаа & Хийх ажил \\
    \midrule
    2026 оны 3-4 сар & AI серверийг MCP хэрэгслүүдтэй бүрэн уях, Chrome extension-ийн token reuse урсгалыг турших, use case-уудыг нарийвчлах \\
    2026 оны 4-5 сар & Ерөнхий MCP өргөтгөлийн загвар, mutation аюулгүй байдлын дүрэм, audit trail логикийг боловсруулах \\
    2026 оны 5 сар & Багшийн порталын API интеграц болон Excel-ээс дүн оруулах workflow-г турших, evaluation scenario бэлтгэх \\
    2026 оны 5-6 сар & Туршилт, хэрэглэгчийн үнэлгээ, алдаа засвар, төгсгөлийн тайлангийн өргөтгөл хийх \\
    \bottomrule
  \end{tabularx}
\end{table}

\section{Үнэлгээний анхны төлөвлөгөө}

Эцсийн шатанд системийг даалгавар амжилттай гүйцэтгэсэн хувь, хариу өгөх хугацаа, хэрэгслийн зөв сонголт, хэрэглэгчийн сэтгэл ханамж, мөн аюулгүй байдлын зөв хэрэгжилт гэсэн үндсэн хэмжүүрээр үнэлэхээр төлөвлөж байна. Өгөгдөл асуух, экспорт хийх, Excel-ээс дүн оруулах зэрэг use case бүр дээр амжилтын хувь, энгийн асуулт болон олон алхамт хүсэлтийн хариу өгөх хугацаа, AI ямар MCP хэрэгсэл ашигласан нь зөв эсэх, оюутнуудын богино асуулга, qualitative feedback, мөн mutation үйлдэл дээр баталгаажуулалт ажилласан эсэх зэргээр үнэлнэ.

Эдгээр хэмжүүр нь дипломын ажлын эцсийн шатанд зөвхөн техникийн demo бус, бодит хэрэглээний үнэ цэнийг нотлоход чухал үзүүлэлт болно.
